Optical and wave-based analog computing has re-emerged as a route to
accelerating linear algebra, with metamaterials and photonic circuits shown to
invert matrices and solve integral equations. However, it has remained unclear
which problems such passive structures can represent, and at what cost, once
encoding, precision, conditioning and readout are accounted for. Here, we show
that every polynomially conditioned sparse linear system can be compiled into a
three-dimensional dielectric -- a cavity filled with lossless, isotropic
material of bounded permittivity -- whose time-harmonic response encodes its
solution. For $m$ nonzero coefficients, the volume of the structure and the
cost of writing it down and reading out the answer scale as $\Ot(m^{3/2})$,
coefficients require only $O(\log m)$ bits, and every field that approximately
satisfies Maxwell's equations decodes to an approximate solution. We prove the
construction for the continuum equations and give a computer-assisted version
for the Yee discretization underlying finite-difference solvers. Read in reverse, it
implies that electromagnetic simulation is, up to this overhead, as hard as
general sparse linear algebra: a sufficiently fast general-purpose Maxwell
solver would yield a faster sparse linear solver.
